\documentclass[11pt,a4paper]{article}

% ============================================================
% Scan-verified LaTeX for Cayley Vol XIII, book pages 56--80
% Source: collmathpapers13caylrich.pdf (PDF pp. 76--100)
% Papers 903 (concl.), 904, 905, 906, 907, 908, 909, 910, 911, 912 (partial)
% ============================================================
\usepackage[utf8]{inputenc}
\usepackage[T1]{fontenc}
\usepackage[english]{babel}
\usepackage{amsmath,amssymb,amsthm}
\usepackage{mathtools}
\usepackage{geometry}
\usepackage{array}
\usepackage{fancyhdr}
\usepackage{hyperref}
\usepackage[protrusion=true,expansion=false]{microtype}

\geometry{
  a4paper,
  left=1in,
  right=1in,
  top=1in,
  bottom=1in,
  footskip=0.5in
}

\hypersetup{
  colorlinks=false,
  pdfborder={0 0 0},
  pdftitle={The Collected Mathematical Papers of Arthur Cayley, Vol. XIII, Pages 56--80},
  pdfauthor={Arthur Cayley}
}

\pagestyle{fancy}
\fancyhf{}
\fancyhead[C]{\small The Collected Mathematical Papers of Arthur Cayley, Vol.\ XIII}
\fancyfoot[C]{\thepage}
\renewcommand{\headrulewidth}{0.4pt}

% Fallback declarations
\providecommand{\sn}{\operatorname{sn}}
\providecommand{\cn}{\operatorname{cn}}
\providecommand{\dn}{\operatorname{dn}}
\providecommand{\AH}{\operatorname{AH}}
\providecommand{\Disct}{\operatorname{Disct.}}

\begin{document}

\thispagestyle{empty}
\begin{center}
\vspace*{2cm}
{\LARGE\bfseries THE COLLECTED MATHEMATICAL PAPERS\\[0.5em] OF ARTHUR CAYLEY, Sc.D., F.R.S.\\[1em]}
{\Large LATE SADLERIAN PROFESSOR OF PURE MATHEMATICS\\ IN THE UNIVERSITY OF CAMBRIDGE.\\[2em]}
{\LARGE\bfseries VOL.\ XIII.\\[2em]}
{\large Pages 56--80\\[0.5em] Papers 903 (conclusion)--912 (partial)\\[2em]}
{\large Scan-verified transcription\\[3em]}
\vfill
{\large Cambridge: At the University Press.\\ 1897.\\[1em]}
\end{center}
\clearpage

\setcounter{page}{56}

% ============================================================
% PAPER 903 (conclusion) -- Pages 56--57
% ON LATIN SQUARES
% ============================================================

\noindent Imagine the square completed: we may write down the substitutions by which we pass from the bottom line to itself (this is of course the substitution 1) and to each of the other lines respectively; we have thus a set of $n$ substitutions, which may form a group; and when this is so, we may conversely from the group construct the Latin square. But it is not every Latin square which is thus connected with a group of $n$ substitutions.

In the cases $n = 2, 3, 4$ there is no difficulty; the squares are
\[
\begin{array}{cccccc}
ba, & cab, & dcab, & dcba, & dcba, & dabc,\\
ab, & bca, & cdba, & cdab, & cadb, & cdab,\\
    & abc, & badc, & badc, & bdac, & bcda,\\
    &      & abcd, & abcd, & abcd, & abcd;
\end{array}
\]
viz.\ when $n = 2$ the number is 1, when $n = 3$ it is 1, when $n = 4$ it is 4; in this last case, the arrangement $badc$ for the second line gives two squares, but each of the other arrangements only one square.

In each of the squares of 4, we have a group, viz.\ for the four squares respectively, these are
\[
\begin{array}{rl}
1^{\circ}, & 1,\ (ab)(cd),\ (acbd),\ (adbc),\\
2^{\circ}, & 1,\ (ab)(cd),\ (ac)(bd),\ (ad)(bc),\\
3^{\circ}, & 1,\ (abdc),\ (acdb),\ (ad)(bc),\\
4^{\circ}, & 1,\ (abcd),\ (ac)(bd),\ (adcb).
\end{array}
\]
$1^{\circ}$, $3^{\circ}$, $4^{\circ}$ are the cyclical groups of $(acbd)$, $(abdc)$, and $(abcd)$, respectively: $2^{\circ}$ is a different kind of group.

In the case when $n = 5$, the whole number of squares is 56; viz.\ there are five arrangements of the second line each giving four squares, and six arrangements each giving six squares, $5 \cdot 4 + 6 \cdot 6 = 56$. The five arrangements are
\[
\begin{array}{ccccc}
baecd, & badec, & bcaed, & bdeac, & bedca,\\
abcde, & abcde, & abcde, & abcde, & abcde,
\end{array}
\]
viz.\ in these cases, the substitutions for passing to the second line are $(ab)(ced)$, $(ab)(cde)$, $(abc)(de)$, $(abd)(ce)$, $(abe)(cd)$, respectively.

The six arrangements are
\[
\begin{array}{cccccc}
bdaec, & beacd, & bcead, & bedac, & bcdea, & bdeca,\\
abcde, & abcde, & abcde, & abcde, & abcde, & abcde;
\end{array}
\]
viz.\ in these cases, the substitutions for passing to the second line are $(abdec)$, $(abedc)$, $(abced)$, $(abecd)$, $(abcde)$, $(abdce)$, respectively.

A set of four squares is
\[
\begin{array}{cccc}
ecdba, & edabc, & ecdab, & edbac,\\
debac, & dcbea, & deabc, & dcaeb,\\
cdaeb, & cedab, & cdbea, & cedba,\\
baecd, & baecd, & baecd, & baecd,\\
abcde, & abcde, & abcde, & abcde;
\end{array}
\]
and a set of six squares is
\[
\begin{array}{cccccc}
ecdab, & eadcb, & ecdba, & eabcd, & eabcd, & ecbad,\\
debca, & dceba, & daecb, & dceba, & dceab, & daecb,\\
caebd, & cebad, & cebad, & cedab, & cedba, & cedba,\\
bdaec, & bdaec, & bdaec, & bdaec, & bdeac, & bdaec,\\
abcde, & abcde, & abcde, & abcde, & abcde, & abcde.
\end{array}
\]

In a square belonging to a set of four, the substitutions for obtaining from any one line all the other lines are of a form such as 1, $(ab)(ced)$, $(ac)(bde)$, $(ad)(bec)$, $(ae)(bcd)$, which are not a group. In the case of a set of six squares, there is one square of the set (in the foregoing instance the first square) where the substitutions are of a form such as 1, $(abdec)$, $(acedb)$, $(adceb)$, $(aebcd)$, and which thus form a cyclical group of five substitutions.

\clearpage

% ============================================================
% PAPER 904 -- Pages 58--59
% NOTE ON RECIPROCAL LINES
% ============================================================
\begin{center}
{\large\bfseries 904.}\\[0.5em]
{\large\bfseries NOTE ON RECIPROCAL LINES.}\\[0.5em]
\small [From the \emph{Messenger of Mathematics}, vol.\ \textsc{xix}.\ (1890), pp.\ 174, 175.]
\end{center}
\medskip

\noindent IF two lines are reciprocal in regard to a quadric surface, then any point on the one line and any point on the other line are harmonics in regard to the surface, viz.\ the two points and the intersections of their line of junction with the surface form a harmonic range. This is obvious: the polar plane of the first point passes through the second line, and thus the second point is a point in the polar plane of the first point, that is, the two points are harmonics.

But it is worth while to look at the theorem in a different point of view. If the first line meets the surface in the points $A$, $C$ and the second line meets the surface in the points $B$, $D$ (in order to the four points being real, the surface must of course be a skew hyperboloid), then $AB$, $BC$, $CD$, $DA$ are lines on the surface, say $AB$, $CD$ are directrices, and $BC$, $AD$ are generatrices, the two reciprocal lines being the diagonals $AC$ and $BD$ of the skew quadrilateral. Taking on $AC$ a point $P$ and on $BD$ a point $Q$, the theorem is that, if $PQ$ meets the surface in the points $X$, $Y$, then the four points $P$, $Q$, $X$, $Y$ are harmonics. Let the generatrix through $X$ meet $AB$, $CD$ in the points $S$, $S'$ respectively, and the generatrix through $Y$ meet the same lines in the points $T$, $T'$ respectively. Then the four lines $CB$, $DA$, $T'T$, $S'S$, are all met by an infinite series of lines, and in particular they are met by the lines $CD$, $BA$ in the points $C$, $D$, $T'$, $S'$ and $B$, $A$, $T$, $S$ respectively. Consequently, writing $\AH$ for anharmonic ratio, we have
\[
\AH(C,\ D,\ T',\ S') = \AH(B,\ A,\ T,\ S).
\]

Consider now the four lines $CA$, $DB$, $T'T$, $S'S$: these are each of them met by the lines $CD$ and $BA$, and also by the line $PQ$: therefore, by an infinity of lines (that is, they are directrices of a hyperboloid). They are met by the last-mentioned three lines in the points $C$, $D$, $T'$, $S'$; $A$, $B$, $T$, $S$; and $P$, $Q$, $Y$, $X$ respectively; hence
\[
\AH(C,\ D,\ T',\ S') = \AH(A,\ B,\ T,\ S) = \AH(P,\ Q,\ Y,\ X).
\]
Comparing with the foregoing equation, it appears that
\[
\AH(B,\ A,\ T,\ S) = \AH(A,\ B,\ T,\ S),
\]
viz.\ the anharmonic ratio is not altered by the interchange of the two points $A$, $B$: this implies that the points $A$, $B$, $T$, $S$ are harmonics, viz.\ the pairs $A$, $B$ and $S$, $T$ are harmonics; and, this being so, the equation
\[
\AH(A,\ B,\ T,\ S) = \AH(P,\ Q,\ Y,\ X)
\]
shows that $P$, $Q$, $Y$, $X$ are harmonics, viz.\ that the pairs $P$, $Q$ and $X$, $Y$ are harmonics; or say $P$, $Q$ are harmonics in regard to the quadric surface, which is the theorem in question.

\clearpage

% ============================================================
% PAPER 905 -- Pages 60--63
% ON THE EQUATION x^{17} - 1 = 0
% ============================================================
\begin{center}
{\large\bfseries 905.}\\[0.5em]
{\large\bfseries ON THE EQUATION $x^{17} - 1 = 0$.}\\[0.5em]
\small [From the \emph{Messenger of Mathematics}, vol.\ \textsc{xix}.\ (1890), pp.\ 184--188.]
\end{center}
\medskip

\noindent WRITING $\rho = \cos\dfrac{2\pi}{17} + i \sin\dfrac{2\pi}{17}$, I carry the solution up to the determination of the periods each of two roots, $\rho + \rho^{16}$, $= 2\cos\dfrac{2\pi}{17}$, \&c. The expressions contain the radicals
\[
a = \sqrt{(17)},\quad b = \sqrt{\{2(17 - a)\}},\quad c = \sqrt{\{4(17 + 3a) - 2(3 + a)b\}},
\]
where $a$, $b$, $c$ are taken to be positive ($a = 4{\cdot}12$, $b = 5{\cdot}07$, $c = 6{\cdot}72$). Taking for a moment $r$ to be any imaginary seventeenth root, $r = \rho^{\theta}$, then the algebraical expression for the period $P_1$ of eight roots is $P_1 = \tfrac{1}{2}(- 1 \pm a)$, but I assume the value to be $P_1 = \tfrac{1}{2}(- 1 + a)$, and thus determine $\theta$ to denote some one of the values 1, 2, 4, 8, 9, 13, 15, 16; similarly, I assume the value of $Q_1$ to be $= \tfrac{1}{4}(- 1 + a + b)$; and thus further determine $\theta$ to denote some one of the values 1, 4, 13, 16: and, again, I assume the value of $R_1$ to be $= \tfrac{1}{8}(- 1 + a + b + c)$, and thus further determine $\theta$ to denote one of the values 1 and 16. As regards the values of the periods $R$, it is obviously indifferent which value is taken, and I assume therefore $\theta = 1$. This comes to saying that the signs of the radicals are determined in suchwise that $r$ shall denote the root $\cos\dfrac{2\pi}{17} + i\sin\dfrac{2\pi}{17}$; and it is to be understood that $r$ has this value.

I write now
\begin{align*}
P_1 &= r + r^9 + r^{13} + r^{15} + r^{16} + r^8 + r^4 + r^2,\\
P_2 &= r^3 + r^{10} + r^5 + r^{11} + r^{14} + r^7 + r^{12} + r^6,\\
Q_1 &= r + r^{13} + r^{16} + r^4,\\
Q_2 &= r^9 + r^{15} + r^8 + r^2,\\
Q_3 &= r^3 + r^5 + r^{14} + r^{12},\\
Q_4 &= r^{10} + r^{11} + r^7 + r^6,\\
R_1 &= r + r^{16},\\
R_2 &= r^{13} + r^4,\\
R_3 &= r^9 + r^8,
\end{align*}
\clearpage
\noindent
\begin{align*}
R_4 &= r^{15} + r^2,\\
R_5 &= r^3 + r^{14},\\
R_6 &= r^5 + r^{12},\\
R_7 &= r^{10} + r^7,\\
R_8 &= r^{11} + r^6,
\end{align*}
and moreover
\begin{align*}
a &= P_1 - P_2,\\
b &= 2(Q_1 - Q_2),\\
+ b_1 &= 2(Q_3 - Q_4),\\
c &= 4(R_1 - R_2),\\
- c_1 &= 4(R_3 - R_4),\\
+ c_2 &= 4(R_5 - R_6),\\
- c_3 &= 4(R_7 - R_8).
\end{align*}

It will appear by what follows that $a$ is determined by the quadric equation $a^2 = 17$, but I have assumed that $a$ denotes the positive root $a = \sqrt{(17)}$; similarly, $b$ is determined by the quadric equation $b^2 = 2(17 - a)$, but it is assumed that $b$ denotes the positive root, $b = \sqrt{\{2(17 - a)\}}$; and $c$ is determined by the quadric equation $c^2 = 4(17 + 3a) - 2(3 + a)b$, but it is assumed that $c$ denotes the positive root,
\[
c = \sqrt{\{4(17 + 3a) - 2(3 + a)b\}}.
\]

If in the equations I had written $b_1$, $c_1$, $c_2$, $c_3$, instead of $+ b_1$, $- c_1$, $+ c_2$, $- c_3$, then $b_1$ comes out rationally in terms of $a$, $b$; and $c_1$, $c_2$, $c_3$ come out rationally in terms of $a$, $b$, $c$; the signs were attached to them \emph{à posteriori}, in suchwise that the values of $b_1$, $c_1$, $c_2$, $c_3$ might be each of them positive; for their independent determination, we have, in fact, for $b_1^2$ an expression such as that for $b^2$; and for $c_1^2$, $c_2^2$, $c_3^2$ expressions such as that for $c^2$; and taking as above for each of them the positive value of the square root, we have
\begin{align*}
a &= \sqrt{(17)},\quad & (&= 4{\cdot}12),\\
b &= \sqrt{\{2(17 - a)\}},\quad & (&= 5{\cdot}07),\\
b_1 &= \sqrt{\{2(17 + a)\}},\quad & (&= 6{\cdot}49),\\
c &= \sqrt{\{4(17 + 3a) - 2(3 + a)b\}},\quad & (&= 6{\cdot}72),\\
c_1 &= \sqrt{\{4(17 + 3a) + 2(3 + a)b\}},\quad & (&= 13{\cdot}77),\\
c_2 &= \sqrt{\{4(17 - 3a) + 2(-3 + a)b_1\}},\quad & (&= 5{\cdot}75),\\
c_3 &= \sqrt{\{4(17 - 3a) - 2(-3 + a)b_1\}},\quad & (&= 2{\cdot}02).
\end{align*}

The relations between the periods $P$ are
\[
P_1 + P_2 = - 1,
\]
\[
\begin{array}{c|c|c}
 & P_1 & P_2 \\\hline
P_1 & -P_1 + 4 & -4 \\\hline
P_2 &          & -P_2 + 4
\end{array}
\qquad
\begin{array}{c|c}
 & a \\\hline
a & 17
\end{array}
\]
that is,
\[
P_1^2 = - P_1 + 4,\ \&\text{c.};\quad a^2 = 17.
\]

Here for the second square, we have
\[
a^2 = (P_1 - P_2)^2 = P_1^2 + P_2^2 - 2 P_1 P_2 = (- P_1 + 4) + (- P_2 + 4) + 8,\ = - P_1 - P_2 + 16,\ = 17;
\]
and similarly in the other cases which follow.

For the periods $Q$, we have
\[
Q_1 + Q_2 = P_1,\quad Q_3 + Q_4 = P_2,
\]
\[
\begin{array}{c|c|c|c|c}
 & Q_1 & Q_2 & Q_3 & Q_4 \\\hline
Q_1 & Q_2 + 2 Q_3 + 4 & -1 & 2 Q_1 + Q_2 + Q_4 & Q_2 + Q_3 + 2 Q_4 \\\hline
Q_2 &  & Q_1 + 2 Q_4 + 4 & Q_1 + 2 Q_3 + 4 & Q_1 + 2 Q_3 + Q_3 \\\hline
Q_3 &  &  & 2 Q_2 + Q_4 + 4 & -1 \\\hline
Q_4 &  &  &  & 2 Q_1 + Q_3 + 4
\end{array}
\]
\[
\begin{array}{c|c|c}
 & b & b_1 \\\hline
b & 34 - 2a & 8a \\\hline
b_1 &  & 34 + 2a
\end{array}
\]
so that $b b_1 = 8a$, or $b_1$ is given rationally in terms of $a$, $b$.

And for the periods $R$, we have
\[
R_1 + R_2 = Q_1,\quad R_3 + R_4 = Q_2,\quad R_5 + R_6 = Q_3,\quad R_7 + R_8 = Q_4,
\]
\[
\renewcommand{\arraystretch}{1.2}
\begin{array}{c|c|c|c|c|c|c|c|c}
 & R_1 & R_2 & R_3 & R_4 & R_5 & R_6 & R_7 & R_8 \\\hline
R_1 & R_4 + 2 & R_3 + R_8 & R_2 + R_7 & R_1 + R_5 & R_2 + R_4 & R_2 + R_8 & R_3 + R_8 & R_8 + R_7 \\\hline
R_2 &  & R_5 + 2 & R_2 + R_8 & R_4 + R_8 & R_1 + R_7 & R_1 + R_3 & R_2 + R_8 & R_4 + R_7 \\\hline
R_3 &  &  & R_1 + 2 & R_7 + R_8 & R_8 + R_7 & R_2 + R_8 & R_1 + R_4 & R_4 + R_5 \\\hline
R_4 &  &  &  & R_2 + 2 & R_1 + R_8 & R_3 + R_7 & R_3 + R_6 & R_2 + R_5 \\\hline
R_5 &  &  &  &  & R_6 + 2 & R_3 + R_4 & R_2 + R_7 & R_8 + R_8 \\\hline
R_6 &  &  &  &  &  & R_7 + 2 & R_4 + R_6 & R_1 + R_8 \\\hline
R_7 &  &  &  &  &  &  & R_8 + 2 & R_1 + R_5 \\\hline
R_8 &  &  &  &  &  &  &  & R_3 + 2
\end{array}
\]
\[
\renewcommand{\arraystretch}{1.4}
\begin{array}{c|c|c|c|c}
 & c & c_1 & c_2 & c_3 \\\hline
c & 4(17 + 3a) - 2(3 + a)b & 8(b + b_1) & 4(2a - b + b_1) & 4(- 2a + b + b_1) \\\hline
c_1 &  & 4(17 + 3a) + 2(3 + a)b & 4(2a + b + b_1) & 4(2a + b - b_1) \\\hline
c_2 &  &  & 4(17 - 3a) + 2(- 3 + a)b_1 & 8(- b + b_1) \\\hline
c_3 &  &  &  & 4(17 - 3a) - 2(- 3 + a)b_1
\end{array}
\]
where observe that the overlined terms $R_5 + R_6$, $R_7 + R_8$, $R_3 + R_4$, and $R_1 + R_2$, have the values $Q_3$, $Q_4$, $Q_2$, $Q_1$, respectively, and in the last table $b_1$ may be considered as denoting its value $= \dfrac{8a}{b}$, so that $c_1$, $c_2$, $c_3$ are each given rationally in terms of $a$, $b$, $c$.

And from the foregoing results, we have
\begin{align*}
P_1 &= \tfrac{1}{2}(- 1 + a), & &= 1{\cdot}56,\\
P_2 &= \tfrac{1}{2}(- 1 - a), & &= - 2{\cdot}56,\\
Q_1 &= \tfrac{1}{4}(- 1 + a + b), & &= 2{\cdot}05,\\
Q_2 &= \tfrac{1}{4}(- 1 + a - b), & &= - 0{\cdot}49,\\
Q_3 &= \tfrac{1}{4}(- 1 - a + b_1), & &= 0{\cdot}34,\\
Q_4 &= \tfrac{1}{4}(- 1 - a - b_1), & &= - 2{\cdot}90,\\
R_1 &= \tfrac{1}{8}(- 1 + a + b + c), & &= 1{\cdot}87,\\
R_2 &= \tfrac{1}{8}(- 1 + a + b - c), & &= 0{\cdot}18,\\
R_3 &= \tfrac{1}{8}(- 1 + a - b - c_1), & &= - 1{\cdot}96,\\
R_4 &= \tfrac{1}{8}(- 1 + a - b + c_1), & &= 1{\cdot}47,\\
R_5 &= \tfrac{1}{8}(- 1 - a + b_1 + c_2), & &= 0{\cdot}89,\\
R_6 &= \tfrac{1}{8}(- 1 - a + b_1 - c_2), & &= - 0{\cdot}55,\\
R_7 &= \tfrac{1}{8}(- 1 - a - b_1 - c_3), & &= - 1{\cdot}70,\\
R_8 &= \tfrac{1}{8}(- 1 - a - b_1 + c_3), & &= - 1{\cdot}20.
\end{align*}

The approximate numerical values have been given throughout only for the purpose of showing that the signs of the square roots have been rightly determined.

\clearpage

% ============================================================
% PAPER 906 -- Pages 64--65
% NOTE ON SCHLAEFLI'S MODULAR EQUATION FOR THE CUBIC TRANSFORMATION
% ============================================================
\begin{center}
{\large\bfseries 906.}\\[0.5em]
{\large\bfseries NOTE ON SCHLAEFLI'S MODULAR EQUATION FOR THE\\ CUBIC TRANSFORMATION; WITH A CORRECTION.}\\[0.5em]
\small [From the \emph{Messenger of Mathematics}, vol.\ \textsc{xx}.\ (1891), pp.\ 59, 60; 120.]
\end{center}
\medskip

\noindent THE equation in question, \emph{Crelle}, t.\ \textsc{lxxii}.\ (1870), p.\ 369, is
\[
S^4 + T^4 + 8 S T - S^3 T^3 = 0,
\]
where
\[
S = \frac{\sqrt{2}}{\sqrt[4]{(k k')}},\quad T = \frac{\sqrt{2}}{\sqrt[4]{(\lambda \lambda')}},
\]
which must be of course equivalent to the ordinary modular equation
\[
u^4 - v^4 + 2 u v - 2 u^3 v^3 = 0,
\]
where
\[
u = \sqrt[4]{k},\quad v = \sqrt[4]{\lambda};
\]
the resemblance in form between the two equations, with such different meanings of the $S$ and $T$ in the one case, and the $u$ and $v$ in the other, is very noticeable.

Schlaefli's equation is
\[
\frac{4}{k k'} + \frac{4}{\lambda \lambda'} + \frac{16}{(k k' \lambda \lambda')^{1/4}} - \frac{8}{(k k' \lambda \lambda')^{3/4}} = 0,
\]
that is,
\[
k k' + \lambda \lambda' = 2 (k k' \lambda \lambda')^{3/4} - 4 (k k' \lambda \lambda')^{1/2},
\]
or say
\[
k^2 k'^2 + \lambda^2 \lambda'^2 = 4 (k k' \lambda \lambda')^{3/2} - 18 (k k' \lambda \lambda').
\]

To deduce this from the $uv$-modular equation, we have (Jacobi's \emph{Fund.\ Nova}, p.\ 68, \emph{Ges.\ Werke}, t.\ \textsc{i}., p.\ 124),
\[
(1 - u^8)(1 - v^8) = (1 - u^2 v^2)^4,
\]
or, multiplying each side by $u^8 v^8$, and extracting the fourth root, we have
\[
\sqrt{(k k' \lambda \lambda')} = u^2 v^2 (1 - u^2 v^2) = x - x^2,
\]
if for shortness we write $x = u^2 v^2$.

The equation to be proved thus is
\[
u^8 (1 - u^8) + v^8 (1 - v^8) = 4 (x - x^2)^3 - 18 (x - x^2)^2 + 16 (x - x^2)^3.
\]
But from the foregoing equation
\[
(1 - u^8)(1 - v^8) = (1 - u^2 v^2)^4,
\]
we have
\[
u^8 + v^8 = 4 u^2 v^2 - 6 u^4 v^4 + 4 u^6 v^6,
\]
and thence
\[
u^{16} + v^{16} = (4 x - 6 x^2 + 4 x^3)^2 - 2 x^4;
\]
and the equation to be proved thus becomes
\[
(4 x - 6 x^2 + 4 x^3) - (4 x - 6 x^2 + 4 x^3)^2 + 2 x^4 = 4 (x - x^2)^3 - 18 (x - x^2)^2 + 16 (x - x^2)^3,
\]
which is in fact an identity, each side being
\[
= 4 x - 22 x^2 + 52 x^3 - 66 x^4 + 48 x^5 - 16 x^6.
\]

\bigskip
\begin{center}
\textsc{Correction, p.\ 120.}
\end{center}

I find that I misquoted Schlaefli's equation, viz.\ in effect, I took it to be
\[
S_1^4 + T_1^4 + 8 S_1 T_1 - S_1^3 T_1^3 = 0,\quad \text{where}\ S_1 = \frac{\sqrt{2}}{\sqrt[4]{(k k')}},\ T_1 = \frac{\sqrt{2}}{\sqrt[4]{(\lambda \lambda')}};
\]
whereas his equation really is
\[
S^4 + T^4 - 8 S T + S^3 T^3 = 0,\quad \text{where}\ S = 2 \sqrt[4]{(k k')},\ T = 2 \sqrt[4]{(\lambda \lambda')}.
\]

The change is only a change of form, for writing $S_1 = \dfrac{2}{S}$ and $T_1 = \dfrac{2}{T}$, the equation in $(S_1, T_1)$ is converted into that in $(S, T)$; but it was quite an unnecessary one, and I cannot account for having made it, as the paper in \emph{Crelle} must have been before me.

\clearpage

% ============================================================
% PAPER 907 -- Page 66
% NOTE ON THE NINTH ROOTS OF UNITY
% ============================================================
\begin{center}
{\large\bfseries 907.}\\[0.5em]
{\large\bfseries NOTE ON THE NINTH ROOTS OF UNITY.}\\[0.5em]
\small [From the \emph{Messenger of Mathematics}, vol.\ \textsc{xx}.\ (1891), p.\ 63.]
\end{center}
\medskip

\noindent LET $\theta$ be a prime ninth root of unity, so that $\theta^6 + \theta^3 + 1 = 0$; and write
\begin{align*}
a &= \theta + \theta^8,\\
b &= \theta^2 + \theta^7,\\
c &= \theta^4 + \theta^5;
\end{align*}
then
\[
a + b + c = \theta + \theta^2 + \cdots + \theta^8 = \frac{1 - \theta^9}{1 - \theta} - 1,\ = - 1,
\]
that is,
\[
a + b + c = 0.
\]
Also
\[
a^2 = b + 2,\quad bc = b - 1,
\]
\[
b^2 = c + 2,\quad ca = c - 1,
\]
\[
c^2 = a + 2,\quad ab = a - 1,
\]
whence
\[
ab + ac + bc = - 3,
\]
\[
abc = - 1;
\]
and $a$, $b$, $c$ are thus the roots of the equation $x^3 - 3 x + 1 = 0$. We have
\[
a^2 b + b^2 c + c^2 a = a^2 + b^2 + c^2 = 6,
\]
\[
a b^2 + b c^2 + c a^2 = bc + ca + ab = - 3,
\]
and thence
\[
- (a^2 b + b^2 c + c^2 a) + (a b^2 + b c^2 + c a^2) = (b - c)(c - a)(a - b) = - 6 - 3,\ = - 9.
\]

The equation $x^3 - 3 x + 1 = 0$ is thus such that $a^2 b + b^2 c + c^2 a$, and consequently any rational function whatever of $a$, $b$, $c$, invariable by the cyclical interchange $(abc)$ of the roots, has a rational value.

\clearpage

% ============================================================
% PAPER 908 -- Pages 67--68
% ON TWO INVARIANTS OF A QUADRIQUADRIC FUNCTION
% ============================================================
\begin{center}
{\large\bfseries 908.}\\[0.5em]
{\large\bfseries ON TWO INVARIANTS OF A QUADRIQUADRIC FUNCTION.}\\[0.5em]
\small [From the \emph{Messenger of Mathematics}, vol.\ \textsc{xx}.\ (1891), pp.\ 68, 69.]
\end{center}
\medskip

\noindent THE quadriquadric function
\begin{align*}
& z^2 \,(a x^2 + 2 h x y + g' y^2)\\
& \quad + 2 z w (h' x^2 + 2 b x y + f y^2)\\
& \quad + w^2 (g x^2 + 2 f' x y + c y^2),
\end{align*}
considered successively as a function of $(z, w)$ and of $(x, y)$, has the discriminants $U$, $V$, equal to
\[
(a x^2 + 2 h x y + g' y^2)(g x^2 + 2 f' x y + c y^2) - (h' x^2 + 2 b x y + f y^2)^2,
\]
\[
(a z^2 + 2 h' z w + g w^2)(g' z^2 + 2 f z w + c w^2) - (h z^2 + 2 b z w + f' w^2)^2,
\]
respectively. As is well known, these quartic functions have each of them the same quadrinvariant and the same cubinvariant; these are the invariants in question of the quadriquadric function.

The quadrinvariant has been calculated in a different notation, but I am not aware that the cubinvariant has been before calculated; the two values are as follows:

\medskip
\noindent\textbf{Quadrinvariant is}
\[
\begin{array}{rr}
a^2 c^2 & +\ 3\\
a b^2 c & -\ 24\\
b^4 & +\ 48\\
a c g g' & +\ 42\\
a c h f' & -\ 12\\
a c h' f & -\ 12\\
a b f f' & +\ 72\\
b c h h' & +\ 72\\
b^2 g g' & -\ 24\\
b^2 f h' & -\ 96\\
b^2 f' h & -\ 96\\
a f^2 g & -\ 36\\
a f'^2 g' & -\ 36\\
b f g h & +\ 72\\
b f g' h' & +\ 72\\
c g h^2 & -\ 36\\
c g' h'^2 & -\ 36\\
g^2 g'^2 & +\ 3\\
f^2 h'^2 & +\ 48\\
f'^2 h^2 & +\ 48\\
g g' h f' & -\ 12\\
g g' h' f & -\ 12\\
f f' h h' & -\ 48\\\hline
\multicolumn{2}{r}{\pm 480}
\end{array}
\]

\medskip
\noindent\textbf{Cubinvariant is}

The cubinvariant has the value $\pm 2866$ obtained by combining the following groups of terms with the indicated coefficients:

\smallskip
\noindent Group 1 ($+76, -49$):
\[
\begin{array}{rr}
a^2 c^3 & -\ 1\\
a^2 b^2 c^2 & +\ 12\\
a b^4 c & -\ 48\\
b^6 & +\ 64
\end{array}
\]

\noindent Group 2 ($+381, -624$):
\[
\begin{array}{rr}
a^2 c^2 g g' & +\ 33\\
a b^2 c g g' & -\ 120\\
a^2 c^2 f h' & +\ 6\\
a^2 c^2 f' h & +\ 6\\
a b^2 c f h' & +\ 24\\
a b^2 c f' h & +\ 24\\
b^4 f h' & -\ 192\\
b^4 f' h & -\ 192\\
b^4 g g' & -\ 48\\
a^2 b c f f' & -\ 36\\
a b c^2 h h' & -\ 36\\
a b^3 f f' & +\ 144\\
b^3 c h h' & +\ 144
\end{array}
\]

\noindent Group 3 ($+648, -432$):
\[
\begin{array}{rr}
a c g^2 g'^2 & +\ 33\\
a c f g' h & +\ 60\\
a c f g h' & +\ 60\\
a c f^2 h'^2 & +\ 24\\
a c f'^2 h^2 & +\ 24\\
a c f f' h h' & +\ 12\\
a b f f' g g' & +\ 180\\
a b f^2 f' h' & -\ 144\\
a b f f'^2 h & -\ 144\\
b c g g' h h' & +\ 180\\
b c f h h'^2 & -\ 144\\
b c f' h^2 h' & -\ 144\\
a^2 f^2 f'^2 & +\ 54\\
c^2 h^2 h'^2 & +\ 54\\
b^2 f^2 h^2 & +\ 192\\
b^2 f'^2 h'^2 & +\ 192\\
b^2 f f' h h' & +\ 96\\
b^2 f g g' h' & +\ 24\\
b^2 f' g g' h & +\ 24
\end{array}
\]

\noindent Group 4 ($+288, -936$):
\[
\begin{array}{rr}
a f^2 f g h & -\ 36\\
a f f'^2 g' h' & -\ 36\\
a f^2 g^2 g' & -\ 36\\
a f g g'^2 & -\ 36\\
a f^3 g h' & +\ 72\\
a f'^3 g' h & +\ 72\\
c f g h^2 h' & -\ 36\\
c f' g' h h'^2 & -\ 36\\
c g^2 g' h^2 & -\ 36\\
c g g'^2 h'^2 & -\ 36\\
c f g h^3 & +\ 72\\
c f' g' h'^3 & +\ 72\\
b f f' g' h'^2 & -\ 144\\
b f^2 g' h h' & -\ 144\\
b f^2 g h h' & -\ 144\\
b f' g h^2 & -\ 144\\
b f g^2 g' h & -\ 36\\
b f' g g'^2 h' & -\ 36
\end{array}
\]

\noindent Group 5 ($+1101, -696$, $+372, -129$):
\[
\begin{array}{rr}
a c f'^2 g & -\ 36\\
a c f^2 g' & -\ 36\\
a c^2 g h^2 & -\ 36\\
a c^2 g' h'^2 & -\ 36\\
a b^2 c f'^2 g & -\ 72\\
a b^2 c f^2 g' & -\ 72\\
a b^2 f^2 g' h' & -\ 144\\
a b^2 f'^2 g h & -\ 144\\
a c f'^2 h^2 & +\ 24\\
a c f^2 h'^2 & +\ 24\\
a f^3 g h' & +\ 72\\
a f'^3 g' h & +\ 72\\
f^4 h'^4 & -\ 64\\
f'^4 h^4 & -\ 64\\
g^3 g'^3 & -\ 1\\
g^2 g'^2 h h' & +\ 6\\
f g^2 g'^2 h' & +\ 6\\
f' g^2 g'^2 h & +\ 6\\
f g g' h'^2 & +\ 24\\
f' g g' h^2 & +\ 24\\
f f' g g' h h' & +\ 12\\
f^3 h^3 & -\ 64\\
f'^3 h'^3 & -\ 64\\
f^2 g h^2 & +\ 54\\
f'^2 g' h'^2 & +\ 54\\
f^2 f' h h'^2 & +\ 96\\
f f'^2 h^2 h' & +\ 96
\end{array}
\]

\medskip
\noindent (The total grand sum is $\pm 2866$, the cubinvariant of the quadriquadric.)

By writing herein $f',\ g',\ h' = f,\ g,\ h$, we obtain of course the two invariants of the symmetrical quadriquadric function.

\clearpage

% ============================================================
% PAPER 909 -- Pages 69--73
% ON A PARTICULAR CASE OF KUMMER'S DIFFERENTIAL EQUATION OF THE THIRD ORDER
% ============================================================
\begin{center}
{\large\bfseries 909.}\\[0.5em]
{\large\bfseries ON A PARTICULAR CASE OF KUMMER'S DIFFERENTIAL\\ EQUATION OF THE THIRD ORDER.}\\[0.5em]
\small [From the \emph{Messenger of Mathematics}, vol.\ \textsc{xx}.\ (1891), pp.\ 75--79.]
\end{center}
\medskip

\noindent THE general form of equation in question is
\[
\frac{x'''}{x'} - \tfrac{3}{2}\left(\frac{x''}{x'}\right)^2 + x'^2\left\{\frac{A}{(x - 1)^2} + \frac{B}{x(x - 1)} + \frac{C}{x^2}\right\} - \left\{\frac{A'}{(t - 1)^2} + \frac{B'}{t(t - 1)} + \frac{C'}{t^2}\right\} = 0,
\]
here $x$ is a function of $t$; and $A$, $B$, $C$, $A'$, $B'$, $C'$ are numerical constants. For various given values of $A$, $B$, $C$, and values determined thereby of $A'$, $B'$, $C'$, the equation admits of a solution in the form $x = $ rational function of $t$; the theory in reference to the cases considered by Schwarz is considered in my paper ``On the Schwarzian Derivative and the Polyhedral Functions,'' \emph{Camb.\ Phil.\ Trans.}, t.\ \textsc{xiii}.\ (1883), pp.\ 5--68, [744]. But the theory is considered in a more general and exhaustive manner in Goursat's memoir, ``Recherches sur l'équation de Kummer,'' \emph{Acta Soc.\ Sci.\ Fennic\ae}, t.\ \textsc{xv}.\ (1888), pp.\ 47--127. I consider here one of the solutions given by him, viz.\ writing
\begin{align*}
P &= 4t - 5,\quad & X &= t^2 P^3,\\
Q &= 5t - 4,\quad & Y &= Q^3,\\
R &= 8t^2 - 11t + 8,\quad & Z &= - (t - 1)^2 R^2,
\end{align*}
so that, identically, $X + Y + Z = 0$; then the solution is expressed by either of the equivalent equations
\[
x = - \frac{X}{Z} = \frac{t^2 P^3}{(t - 1)^2 R^2},
\]
\[
x - 1 = \frac{Y}{Z} = - \frac{Q^3}{(t - 1)^2 R^2}.
\]

The values of the constants to which this solution belongs are
\[
A = \tfrac{4}{9},\quad B = - \tfrac{32}{63},\quad C = \tfrac{4}{9};\quad A' = \tfrac{4}{9},\quad B' = \tfrac{131}{144},\quad C' = \tfrac{5}{9}.
\]

But instead of assuming these values in the first instance, I leave the values indeterminate; and starting from the foregoing expression for $x$, I substitute this in
\[
\Omega = \frac{x'''}{x'} - \tfrac{3}{2}\!\left(\frac{x''}{x'}\right)^2 + x'^2\!\left\{\frac{A}{(x - 1)^2} + \frac{B}{x(x - 1)} + \frac{C}{x^2}\right\} - \left\{\frac{A'}{(t - 1)^2} + \frac{B'}{t(t - 1)} + \frac{C'}{t^2}\right\},
\]
thus obtaining $\Omega$ as a function of $t$ which, as will appear, vanishes identically when $A$, $B$, $C$, $A'$, $B'$, $C'$ have the foregoing values.

I remark that this is, in effect, doing in a somewhat different form for the particular case what Goursat does for the general case, viz.\ starting from
\[
\Omega_1 = \frac{x'''}{x'} - \tfrac{3}{2}\!\left(\frac{x''}{x'}\right)^2 + x'^2\!\left\{\frac{A}{(x - 1)^2} + \frac{B}{x(x - 1)} + \frac{C}{x^2}\right\},
\]
with values of $A$, $B$, $C$ which belong to the solution considered, he shows that this is a function of $t$ having no infinities other than $(0, 1, \infty)$; that $\infty$ is not an infinity of the function or of the function multiplied into $t$, and that 0 and 1 are each of them a twofold infinity; that is, that the function is of the form
\[
\frac{L t^2 + M t + N}{t^2 (t - 1)^2}\ \text{or}\ \frac{A'}{(t - 1)^2} + \frac{B'}{t(t - 1)} + \frac{C'}{t^2}.
\]

Proceeding to carry out the process, we have
\[
\frac{x'}{x} = \frac{2}{t} + \frac{3 P'}{P} - \frac{2 R'}{R},
\]
\[
\frac{x'}{x - 1} = - \frac{1}{t - 1} + \frac{3 Q'}{Q} - \frac{2 R'}{R},
\]
\[
\frac{x'}{x} = \frac{2}{t} - \frac{1}{t - 1} + \frac{3 P'}{P} - \frac{2 R}{R},
\]
and from either of these equations, collecting and reducing,
\[
x' = \frac{5 t^2 P^2 Q^2}{(t - 1)^3 R^3},
\]
where observe that, from the values of $x$ and $x - 1$ respectively, it appears \emph{à priori} that $t P^2$ and $Q^2$ must be factors in the numerator of $x'$. From this value of $x'$, we have
\[
\frac{x''}{x'} = \frac{1}{t} - \frac{2}{t - 1} + \frac{2 P'}{P} + \frac{2 Q'}{Q} - \frac{3 R'}{R};
\]
and hence, $P'$ and $Q'$ being mere constants,
\[
\frac{x'''}{x'} - \!\left(\frac{x''}{x'}\right)^2 = - \frac{1}{t^2} + \frac{2}{(t - 1)^2} - \frac{3 R''}{R} - 2\!\left(\frac{P'}{P}\right)^2 - 2\!\left(\frac{Q'}{Q}\right)^2 + 3\!\left(\frac{R'}{R}\right)^2,
\]
and consequently
\[
\frac{x'''}{x'} - \tfrac{3}{2}\!\left(\frac{x''}{x'}\right)^2 = - \frac{1}{t^2} + \frac{2}{(t - 1)^2} - \frac{3 R''}{R} - 2\!\left(\frac{P'}{P}\right)^2 - 2\!\left(\frac{Q'}{Q}\right)^2 + 3\!\left(\frac{R'}{R}\right)^2
\]
\[
\hspace{2cm} - \tfrac{1}{2}\!\left(\frac{1}{t} - \frac{2}{t - 1} + \frac{2 P'}{P} + \frac{2 Q'}{Q} - \frac{3 R'}{R}\right)^2.
\]
Putting for shortness
\[
\frac{1}{t} = \alpha,\quad \frac{1}{t - 1} = \beta,\quad \frac{P'}{P} = p,\quad \frac{Q'}{Q} = q,\quad \frac{R'}{R} = r,
\]
this equation gives
\begin{align*}
\Omega = & - \alpha^2 + 2 \beta^2 - \frac{3 R''}{R} - 2 p^2 - 2 q^2 + 3 r^2\\
& - \tfrac{1}{2}(\alpha - 2 \beta + 2 p + 2 q - 3 r)^2\\
& + A(- \beta + 3 q - 2 r)^2\\
& + B(- \beta + 3 q - 2 r)(2 \alpha - \beta + 3 p - 2 r)\\
& + C(2 \alpha - \beta + 3 p - 2 r)^2\\
& - C' \alpha^2 - B' \alpha \beta - A' \beta^2,
\end{align*}
which is
\begin{align*}
= & \ \alpha^2\!\left(- \tfrac{1}{2} + 4 C - C'\right) + \frac{Lt^2}{R}\ \text{say it is}\ = L \alpha^2 - \frac{48}{R}\\
& + \alpha \beta(2 - 2 B - 4 C - B') & & + M \alpha \beta\\
& + \beta^2(A + B + C - A') & & + N \beta^2\\
& + \alpha p(- 2 + 12 C) & & + F \alpha p\\
& + \alpha q(- 2 + 6 B) & & + G \alpha q\\
& + \alpha r(3 - 4 B - 8 C) & & + H \alpha r\\
& + \beta p(4 - 3 B - 6 C) & & + F' \beta p\\
& + \beta q(4 - 6 A - 3 B) & & + G' \beta q\\
& + p^2(- 4 + 9 C) & & + A'' p^2\\
& + q^2(- 4 + 9 A) & & + B'' q^2\\
& + r^2(- 1 + 4 A + 4 B + 4 C) & & + C'' r^2\\
& + q r(6 - 12 A - 6 B) & & + F'' q r\\
& + r p(6 - 6 B - 12 C) & & + G'' r p
\end{align*}

By decomposing $\alpha \beta$, $\alpha p$, \&c., into simple fractions, this becomes
\begin{align*}
\Omega = &\ \ Lx - \frac{48}{R}\\
&+ M(- \alpha + \beta)\\
&+ N \beta^2\\
&+ F(- \tfrac{5}{4} \alpha + \tfrac{5}{4} p)\\
&+ G(- \tfrac{4}{5} \alpha + \tfrac{4}{5} q)\\
&+ H\!\left(- \tfrac{11}{8} \alpha + \frac{11 t + \tfrac{5}{8}}{R}\right)\\
&+ F'(- 4 \beta + 4 p)\\
&+ G'\,(\phantom{-}5 \beta - 5 q)\\
&+ H'\!\left(\phantom{-}\beta + \frac{- 8 t + 19}{R}\right)\\
&+ A'' p^2\\
&+ B'' q^2\\
&+ C'' r^2\\
&+ F''\,\tfrac{8}{5}\!\left(q - \frac{8 t - 43}{R}\right)\\
&+ G''\,\tfrac{8}{5}\!\left(p - \frac{8 t - 13}{R}\right)\\
&+ H''\,\tfrac{8 t}{5} (p - q).
\end{align*}
This is
\begin{align*}
= &\ \ \alpha L\\
&+ \alpha(- M - \tfrac{5}{4} F - \tfrac{4}{5} G - \tfrac{11}{8} H)\\
&+ \beta^2 N\\
&+ \beta\,(M - 4 F' + 5 G' + H')\\
&+ p^2 A''\\
&+ p(4 F' + \tfrac{5}{4} G'' + \tfrac{8 t}{5} H'' + \tfrac{8}{5} F)\\
&+ q^2 B''\\
&+ q(- 5 G' - \tfrac{8}{5} F'' - \tfrac{8 t}{5} H'' + \tfrac{4}{5} G)\\
&+ r^2 C''\\
&+ \frac{1}{R}\!\left[- 48 + H(11 t + \tfrac{5}{8}) + H'(- 8 t + 19) - \tfrac{8}{5} F''(8 t - 43) - \tfrac{8}{5} G''(8 t - 13)\right].
\end{align*}
This should be identically $= 0$; making $A'' = 0$, $B'' = 0$, $C'' = 0$, we find
\[
A = \tfrac{4}{9},\quad B = - \tfrac{32}{63},\quad C = \tfrac{4}{9};\quad (A + B + C = \tfrac{8}{63}).
\]
and thence
\[
F = \tfrac{4 \cdot 9}{9},\quad G = - \tfrac{2 \cdot 4}{63},\quad H = \tfrac{2}{3};\quad F' = \tfrac{2 \cdot 9}{63},\quad G' = - \tfrac{1 \cdot 2}{63},\quad H' = - \tfrac{1}{3};\quad F'' = \tfrac{1 \cdot 2}{63},\quad G'' = \tfrac{1 \cdot 2}{63},\quad H'' = - \tfrac{2}{3}.
\]

These values make the coefficients of $p$ and $q$ to be each $= 0$; and they make the coefficient of $R$ to be identically $= 0$, viz.\ we have
\[
0 = - 48 + \tfrac{1}{8} H + 19 H' + \tfrac{344}{5} F'' + \tfrac{104}{5} G'',
\]
and
\[
0 = \phantom{-}11 H - 8 H' - \tfrac{64}{5} F'' - \tfrac{64}{5} G''.
\]
We have, moreover,
\[
L = \tfrac{1}{2} - C',\quad M = \tfrac{131}{144} - B',\quad N = \tfrac{4}{9} - A';
\]
and the coefficients of $\alpha$ and $\beta$ are $= - M + \tfrac{131}{144}$ and $M - \tfrac{131}{144}$ respectively; hence the coefficients of $\alpha^2$, $\beta^2$, $\alpha$ and $\beta$ will all vanish if only $L = 0$, $M = \tfrac{131}{144}$, $N = 0$, that is,
\[
A' = \tfrac{4}{9},\quad B' = \tfrac{131}{144},\quad C' = \tfrac{5}{9};
\]
and we have thus identically $\Omega = 0$, if only $A$, $B$, $C$, $A'$, $B'$, $C'$ have the above-mentioned values.

\clearpage

% ============================================================
% PAPER 910 -- Pages 74--75
% NOTE ON THE INVOLUTANT OF TWO BINARY MATRICES
% ============================================================
\begin{center}
{\large\bfseries 910.}\\[0.5em]
{\large\bfseries NOTE ON THE INVOLUTANT OF TWO BINARY MATRICES.}\\[0.5em]
\small [From the \emph{Messenger of Mathematics}, vol.\ \textsc{xx}.\ (1891), pp.\ 136, 137.]
\end{center}
\medskip

\noindent CONSIDER the two matrices
\[
M = \begin{pmatrix} a, & b \\ c, & d \end{pmatrix},\quad M' = \begin{pmatrix} a', & b' \\ c', & d' \end{pmatrix},
\]
and their product in one or the other order
\[
M M' = \begin{pmatrix} A, & B \\ C, & D \end{pmatrix},\quad M' M = \begin{pmatrix} A_1, & B_1 \\ C_1, & D_1 \end{pmatrix}.
\]
Then the Involutant is by definition $=$ either of the determinants
\[
I = \begin{vmatrix} 1, & a, & a', & A \\ 0, & b, & b', & B \\ 0, & c, & c', & C \\ 1, & d, & d', & D \end{vmatrix},\quad I_1 = \begin{vmatrix} 1, & a', & a, & A_1 \\ 0, & b', & b, & B_1 \\ 0, & c', & c, & C_1 \\ 1, & d', & d, & D_1 \end{vmatrix};
\]
viz.\ it is to be shown that these two values are in fact equal.

We have
\[
M M' = \begin{pmatrix} a, & b \\ c, & d \end{pmatrix} \begin{pmatrix} a', & b' \\ c', & d' \end{pmatrix} = \begin{pmatrix} A, & B \\ C, & D \end{pmatrix},
\]
that is,
\[
A = a a' + b c',\quad B = a b' + b d',
\]
\[
C = c a' + d c',\quad D = c b' + d d',
\]
and similarly
\[
M' M = \begin{pmatrix} a', & b' \\ c', & d' \end{pmatrix} \begin{pmatrix} a, & b \\ c, & d \end{pmatrix} = \begin{pmatrix} A_1, & B_1 \\ C_1, & D_1 \end{pmatrix},
\]
that is,
\[
A_1 = a a' + c b',\quad B_1 = b a' + d b',
\]
\[
C_1 = a c' + c d',\quad D_1 = b c' + d d',
\]
viz.\ $A_1$, $B_1$, $C_1$, $D_1$ are obtained from $A$, $B$, $C$, $D$ by the interchange of the accented and unaccented letters.

We have then, from the first expression for the Involutant,
\begin{align*}
I &= A \begin{vmatrix} 0, & b, & b' \\ 0, & c, & c' \\ 1, & d, & d' \end{vmatrix} - B \begin{vmatrix} 0, & c, & c' \\ 1, & d, & d' \\ 0, & b, & b' \end{vmatrix} + C \begin{vmatrix} 1, & d, & d' \\ 0, & b, & b' \\ 0, & c, & c' \end{vmatrix} - D \begin{vmatrix} 1, & a, & a' \\ 0, & b, & b' \\ 0, & c, & c' \end{vmatrix}\\
&= A (bc' - b'c) - B (ac' - a'c + cd' - c'd) + C (ab' - a'b + bd' - b'd) - D (bc' - b'c);
\end{align*}
or substituting for $A - D$, $B$ and $C$ their values, this is
\begin{align*}
&= (a a' - d d' + b c' - b' c)(b c' - b' c) - (a b' + b d')(a c' - a' c + c d' - c' d)\\
&\quad + (c a' + d c')(a b' - a' b + b d' - b' d);
\end{align*}
and multiplying out and grouping together the terms in $bc$, $b'c'$, $bc'$ and $b'c$, this is found to be
\[
= - (a' - d')^2 b c + (a - d)(a' - d')(b c' + b' c) - (a - d)^2 b' c' + (b c' - b' c)^2,
\]
which is
\[
= - \{(a - d) b' - (a' - d') b\}\{(a - d) c - (a' - d') c'\} + (b c' - b' c)^2.
\]

Hence, writing
\[
a = b c' - b' c,\quad f = a d' - a' d,
\]
\[
b = c a' - c' a,\quad g = b d' - b' d,
\]
\[
c = a b' - a' b,\quad h = c d' - c' d,
\]
we have
\[
I = -(c + g)(- b + h) + a^2,
\]
that is,
\[
I = bc - ch + bg - gh + a^2.
\]

To obtain the value of $I_1$, we must interchange the accented and unaccented letters, that is, change the signs of the several quantities $a$, $b$, $c$, $f$, $g$, $h$; but $I$, being a quadric function of the six quantities, is not altered by the change; that is, we have $I = I_1$.

\clearpage

% ============================================================
% PAPER 911 -- Pages 76--78
% ON AN ALGEBRAICAL IDENTITY RELATING TO THE SIX COORDINATES OF A LINE
% ============================================================
\begin{center}
{\large\bfseries 911.}\\[0.5em]
{\large\bfseries ON AN ALGEBRAICAL IDENTITY RELATING TO THE SIX\\ COORDINATES OF A LINE.}\\[0.5em]
\small [From the \emph{Messenger of Mathematics}, vol.\ \textsc{xx}.\ (1891), pp.\ 138--140.]
\end{center}
\medskip

\noindent THE following identity may be verified without difficulty; but it is interesting in regard to the analytical theory of the line.

The identity is
\begin{align*}
0 = &\ \begin{vmatrix} B, & C, & F \\ b, & c, & f \\ b', & c', & f' \end{vmatrix} \begin{vmatrix} A, & B, & C \\ a, & b, & c \\ a', & b', & c' \end{vmatrix}\\
&- (bc' - b'c)^2 (AF + BG + CH)\\
&+ (bc' - b'c)(bC - cB)(Af + Bg' + Ch' + Fa' + Gb' + Hc')\\
&- (bc' - b'c)(b'C - c'B)(Af + Bg + Ch + Fa + Gb + Hc)\\
&- (bC - cB)^2 (a'f + b'g' + c'h')\\
&+ (bC - cB)(b'C - c'B)(af + a'f' + bg' + b'g + ch' + c'h)\\
&- (b'C - c'B)^2 (af + bg + ch),
\end{align*}
where all the letters have arbitrary values.

It follows that, if
\[
af + bg + ch = 0,
\]
\[
a'f' + b'g' + c'h' = 0,
\]
\[
af' + a'f + bg' + b'g + ch' + c'h = 0,
\]
\[
AF + BG + CH = 0,
\]
\[
Af + Bg + Ch + Fa + Gb + Hc = 0,
\]
\[
Af' + Bg' + Ch' + Fa' + Gb' + Hc' = 0,
\]
then either
\[
\begin{vmatrix} B, & C, & F \\ b, & c, & f \\ b', & c', & f' \end{vmatrix} = 0,
\]
or else
\[
\begin{vmatrix} A, & B, & C \\ a, & b, & c \\ a', & b', & c' \end{vmatrix} = 0.
\]

Supposing the first three of the six equations are satisfied, then $(a, b, c, f, g, h)$ and $(a', b', c', f', g', h')$ are the coordinates of two intersecting lines; and supposing that the last three equations are also satisfied, then $(A, B, C, F, G, H)$ will be the coordinates of a line meeting each of the intersecting lines. The third line is thus either in the plane of the intersecting lines, or else passes through their point of intersection; and in fact, the first of the two determinant equations is the condition in order that the line may be in the plane of the two intersecting lines, and the second determinant equation is the condition in order that it may pass through their point of intersection. Each equation is in fact one out of four equivalent forms, viz.\ we may have in the first equation $(B, C, F)$, $(C, A, G)$, $(A, B, H)$, or $(F, G, H)$; and in the second equation $(A, B, C)$, $(A, H, G)$, $(H, B, F)$, or $(G, F, C)$.

The analytical theory may be presented in a complete form by means of the formulae of my memoir ``On the six coordinates of a line,'' (1867), \emph{Camb.\ Phil.\ Trans.}, t.\ \textsc{xi}.\ pp.\ 290--323, [435]. Considering the two intersecting lines $(a, b, c, f, g, h)$ and $(a', b', c', f', g', h')$, the coordinates of the plane through these two lines (that is, the coefficients $\xi, \eta, \zeta, \omega$ of the equation $\xi x + \eta y + \zeta z + \omega w = 0$ of the plane) are there given (see p.\ 295\,*) in a fourfold form; and if we thence form the condition in order that the line $(A, B, C, F, G, H)$ may lie in this plane, we have
\[
\begin{vmatrix} 0, & C, & -B, & F \\ -C, & 0, & A, & G \\ B, & -A, & 0, & H \\ -F, & -G, & -H, & 0 \end{vmatrix}
\]
\[
\left(\begin{array}{cccc}
af' + b'g + c'h, & bf - b'f', & cf - c'f', & -(bc' - b'c) \\
ag' - a'g, & a'f + bg' + c'h, & cg' - c'g, & -(ca' - c'a) \\
ah' - a'h, & bh' - b'h, & a'f + b'g + ch', & -(ab' - a'b) \\
gh' - g'h, & hf - h'f', & fg' - f'g, & af + bg' + ch'
\end{array}\right) = 0;
\]
viz.\ the condition is expressible in any one of the 16 forms obtained by combining a line of the first matrix with a line of the second matrix; thus one form is
\[
0 \cdot (af' + b'g + c'h) + C(bf - b'f') - B(cf - c'f') - F(bc' - b'c) = 0,
\]
that is,
\[
\begin{vmatrix} B, & C, & F \\ b, & c, & f \\ b', & c', & f' \end{vmatrix} = 0,
\]
the foregoing first determinant equation, which thus belongs to the case where the line lies in the plane of the two intersecting lines.

\smallskip
\noindent{\footnotesize [* This Collection, vol.\ \textsc{vii}., p.\ 71.]}

\medskip
Again we have (see p.\ 296\,*) an expression, in a fourfold form, for the coordinates of the point of intersection of the two intersecting lines; and thence for the condition, in order that the line $(A, B, C, F, G, H)$ may pass through this point, we have
\[
\begin{vmatrix} 0, & H, & -G, & A \\ -H, & 0, & F, & B \\ G, & -F, & 0, & C \\ -A, & -B, & -C, & 0 \end{vmatrix}
\]
\[
\left(\begin{array}{cccc}
af' + b'g + c'h, & ag' - a'g, & ah' - a'h, & -(bc' - b'c) \\
bf - b'f', & a'f + bg' + c'h, & bh' - b'h, & -(ca' - c'a) \\
cf - c'f', & cg' - c'g, & a'f + b'g + ch', & -(ab' - a'b) \\
gh' - g'h, & hf - h'f', & fg' - f'g, & af + bg' + ch'
\end{array}\right) = 0,
\]
viz.\ the condition is expressible in any one of the 16 forms obtained by combining a line of the first matrix with a line of the second matrix; thus one form is
\[
A(bc' - b'c) + B(ca' - c'a) + C(ab' - a'b) + 0 \cdot (af + bg' + ch') = 0,
\]
that is,
\[
\begin{vmatrix} A, & B, & C \\ a, & b, & c \\ a', & b', & c' \end{vmatrix} = 0,
\]
the foregoing second determinant equation, which thus belongs to the case where the line passes through the point of intersection of the two intersecting lines.

I remark that the original identity may be written in the very compendious symbolical form
\[
\begin{vmatrix} A, & a, & a' \\ F, & f, & f' \end{vmatrix} \begin{vmatrix} B, & b, & b' \\ C, & c, & c' \end{vmatrix} = \begin{vmatrix} A, & a, & a' \\ B, & b, & b' \\ C, & c, & c' \end{vmatrix} \begin{vmatrix} F, & f, & f' \\ B, & b, & b' \\ C, & c, & c' \end{vmatrix}
\]
viz.\ here on the left-hand side the second factor denotes the three determinants $bc' - b'c$, $b'C - c'B$, $Bc - Cb$: the whole is a quadric function of these, the coefficients being
\[
AF + BG + CH,\quad af + bg + ch,
\]
\[
a'f' + b'g' + c'h',\quad af' + bg' + ch' + f'a + g'b + h'c,
\]
\[
a'F + b'G + c'H + f'A + g'B + h'C,
\]
\[
Af + Bg + Ch + Fa + Gb + Hc,
\]
respectively; and the right-hand side is simply the product of two determinants.

\smallskip
\noindent{\footnotesize [* Loc.\ cit., p.\ 72.]}

\clearpage

% ============================================================
% PAPER 912 -- Pages 79--80 (partial, continues beyond p. 100)
% ON THE NOTION OF A PLANE CURVE OF A GIVEN ORDER
% ============================================================
\begin{center}
{\large\bfseries 912.}\\[0.5em]
{\large\bfseries ON THE NOTION OF A PLANE CURVE OF A GIVEN ORDER.}\\[0.5em]
\small [From the \emph{Messenger of Mathematics}, vol.\ \textsc{xx}.\ (1891), pp.\ 148--150.]
\end{center}
\medskip

\noindent WE have a complete geometrical notion of a curve of a given order, viz.\ a curve of the order $n$ is a curve which is met by any line whatever in $n$ points and no more; but starting with this definition, how do we know that there exists a curve of the order $n$? and, further, how do we know that it depends linearly on $\tfrac{1}{2} n (n + 3)$ parameters, or, what is the same thing, that there is one and only one curve which can be drawn through $\tfrac{1}{2} n (n + 3)$ given points?

The last-mentioned property does not by itself constitute a definition of a curve of the order $n$; thus $n = 2$, we cannot define a curve of the second order as a curve which is uniquely determined by the condition of passing through 5 given points; for a cubic passing through 4 given points is a curve uniquely determined by the condition in question; but we may differentiate between these two solutions by adding the further condition that, when 3 of the 5 points are in a line, the curve of the second order shall include as part of itself this line. And we are thus led to the definition: A curve of the order $n$ is a curve which is uniquely determined by the condition of passing through $\tfrac{1}{2} n (n + 3)$ given points; and of being moreover such that, when $n + 1$ of these points lie on a line, it includes as part of itself this line.

Starting from the foregoing definition, the first property is, I think, demonstrable, viz.\ the property that a curve of the order $n$ is met by any line whatever in $n$ points and no more. Thus $n = 2$: start with a line chosen at pleasure, and on it take 2 points which are regarded as indeterminate points: to fix the ideas, let one of these be regarded as depending on a parameter $\lambda$, and the other on a parameter $\mu$, (so that when $\lambda$ has a determinate value assigned to it, the first point becomes a determinate point, and similarly when $\mu$ has a determinate value assigned to it, the second point becomes a determinate point; and consequently, when $\lambda$ and $\mu$ have determinate values, the 2 points are determinate points on the line). Take now any other 3 points not on the line; then, for the moment regarding the 2 points on the line as determinate points, we can through the 5 points draw a curve of the second order; this is the general curve of the second order through the 3 points; for it is a curve of the second order through the 3 points, and which, when the parameters $\lambda$ and $\mu$ are regarded as undetermined and arbitrary, or choosable at pleasure, might be made to pass through any other two points whatever. But by hypothesis, the curve meets the line in question, that is, any line, in two points: and the conic cannot meet the line in more than two points; for in like manner, starting with the given line, and upon it 3 points (which may be considered as depending on the parameters $\lambda$, $\mu$ and $\nu$ respectively), and taking any two points not on the line, we have through the 5 points a conic; and this conic, regarding the parameters $\lambda$, $\mu$, $\nu$ as undetermined and arbitrary, or choosable at pleasure, will be the general conic through the two points; for by a proper determination of the parameters, it might be made to pass through any other 3 points whatever. But, by hypothesis, the conic contains as part of itself the line; that is, the general conic through the 2 points contains as part of itself any line whatever, which is absurd.

So again for the cubic, $n = 3$: here starting with a line taken at pleasure, we take on it 3 points, which may be regarded as depending on the parameters $\lambda$, $\mu$, $\nu$ respectively, and we take any other 6 points not on the line. We have through the 9 points a cubic; and this is the general cubic through the 6 points, for it depends on the parameters $\lambda$, $\mu$, $\nu$, which might be determined so as to make the curve pass through any other 3 points whatever; and by hypothesis the cubic meets the line, that is, any line whatever, in 3 points. And it cannot meet it in more than 3 points: for starting with the same line and upon it 4 points depending on the parameters $\lambda$, $\mu$, $\nu$, $\rho$ respectively, and taking any other 5 points not upon the line, we then have through the 9 points a cubic which will be the general cubic through the 5 points (for it depends on the 4 parameters $\lambda$, $\mu$, $\nu$, $\rho$). But by hypothesis, the cubic contains as part of itself the line, viz.\ the general cubic through the 5 points contains as part of itself any line whatever, which is absurd. The reasoning is quite general; and applying to a curve of the order $n$, the conclusion is that such a curve meets any line in $n$ points and no more.

\end{document}
